\subsection{The sharp norm of the signed dilation operator}
\label{ssec:renewal-norm}

The coefficients in \eqref{eq:lambda-log-renewal} can be grouped by
prime powers. The exact identity is
\begin{equation}\label{eq:prime-renewal-full}
 \sum_{n\ge1}\lambda(n)\log n\,e^{-n/X}
 =\sum_{p^k\ge2}(-1)^k\log p\,B(X/p^k).
\end{equation}
Indeed $\lambda\log=\lambda*(\lambda\Lambda)$, and
$\lambda(p^k)\Lambda(p^k)=(-1)^k\log p$. Absolute convergence
at fixed $X$ follows by comparison with
$\sum_n n\log(2n)e^{-n/X}$.

Consider a finite truncation on a weighted Hilbert space:
\[
 \mathcal K_\sigma=L^2((0,\infty),X^{-2\sigma}\dd X/X),\qquad
 T_R f(X)=\sum_{p^k\le R}(-1)^k\log p\,f(X/p^k).
\]
The following calculation retains every signed prime-power term.

\begin{theorem}[Exact dilation norm]\label{thm:sharp-dilation-norm}
For every real $\sigma$ and finite $R\ge2$,
\begin{equation}\label{eq:sharp-dilation-norm}
 \|T_R\|_{\mathcal K_\sigma\to\mathcal K_\sigma}
       =\sum_{p^k\le R}(\log p)p^{-k\sigma}.
\end{equation}
For each fixed $1/2<\sigma<1$, as $R\to\infty$,
\begin{equation}\label{eq:dilation-norm-growth}
 \|T_R\|\sim\frac{R^{1-\sigma}}{1-\sigma},\qquad
 \frac{\|T_R\|}{\log R}\longrightarrow\infty.
\end{equation}
\end{theorem}
\begin{proof}
The map $Uf(t)=e^{-\sigma t}f(e^t)$ is an isometry onto
$L^2(\mathbb R)$. Under it, $f(X/d)$ becomes
$d^{-\sigma}(Uf)(t-\log d)$. With the Fourier convention using
$e^{-it\tau}$, the multiplier of $T_R$ is
\[
 m_{R,\sigma}(\tau)=\sum_{p^k\le R}
 (-1)^k(\log p)p^{-k\sigma}e^{-ik\tau\log p}.
\]
Plancherel identifies the operator norm with its essential supremum;
see \cite[\S1.14]{dlmf-transforms} for the transform conventions.
The triangle inequality proves the upper bound in
\eqref{eq:sharp-dilation-norm}.

The logarithms of the finitely many primes at most $R$ are linearly
independent over $\mathbb Q$, by unique factorization. Thus the
continuous phase orbit $(e^{-i\tau\log p})_p$ is dense in the finite
product of unit circles. Here is a direct justification. Every
nonconstant integer Fourier character of this orbit has the form
$e^{-i\tau\sum_p m_p\log p}$ with a nonzero frequency. Its average
over $[-T,T]$ tends to zero. Trigonometric-polynomial approximation
therefore gives the Haar average for each continuous function on the
torus. A nonnegative continuous bump supported in any nonempty open
set has positive Haar integral, so the orbit meets that set.

In particular, all prime phases can approach $-1$ simultaneously.
Then $(-1)^ke^{-ik\tau\log p}$ approaches $1$ for every term of
the finite sum. The multiplier approaches the positive sum in
\eqref{eq:sharp-dilation-norm}. Its continuity implies the same lower
bound for the essential supremum, proving equality.

For the asymptotic, the prime number theorem
\cite[\S27.12]{dlmf-primes} and partial summation give
\[
 \sum_{p\le R}(\log p)p^{-\sigma}
       \sim\frac{R^{1-\sigma}}{1-\sigma}.
\]
The higher prime powers contribute a bounded amount because
\[
 \sum_p\sum_{k\ge2}(\log p)p^{-k\sigma}
 \le\frac1{1-2^{-\sigma}}
             \sum_{n\ge2}(\log n)n^{-2\sigma}<\infty.
\]
This proves \eqref{eq:dilation-norm-growth}.
\end{proof}

The size in this theorem is not confined to very high logarithmic
frequencies.

\begin{proposition}[A low-frequency sign change]
\label{prop:renewal-low-frequency}
Fix $1/2<\sigma<1$, and put
\[
 P_R=\sum_{p\le R}(\log p)p^{-\sigma},\qquad
 \tau_R=\frac\pi{\log R}.
\]
Then
\begin{align}
 m_{R,\sigma}(\tau_R)
   &=P_R+O_\sigma\!\left(\frac{R^{1-\sigma}}{\log R}\right),
       \label{eq:renewal-small-frequency}\\
 m_{R,\sigma}(0)&=-P_R+O_\sigma(1).
       \label{eq:renewal-zero-frequency}
\end{align}
If $\psi$ is a fixed continuous smoothing multiplier near zero with
$\psi(0)\ne0$, then
$|\psi(\tau_R)m_{R,\sigma}(\tau_R)|\sim|\psi(0)|P_R$.
\end{proposition}
\begin{proof}
The prime part of $m_{R,\sigma}(\tau_R)$ is
\[
 -\sum_{p\le R}(\log p)p^{-\sigma}e^{-i\tau_R\log p}
 =\sum_{p\le R}(\log p)p^{-\sigma}
           e^{i\pi\log(R/p)/\log R}.
\]
Its difference from $P_R$ has absolute value at most
\[
 \frac\pi{\log R}\sum_{p\le R}
             (\log p)p^{-\sigma}\log(R/p).
\]
Chebyshev's upper bound and partial summation give
$P_u\ll_\sigma u^{1-\sigma}$. Interchanging a finite sum and an
integral consequently gives
\[
 \sum_{p\le R}(\log p)p^{-\sigma}\log(R/p)
   =\int_2^R P_u\frac{\dd u}u\ll_\sigma R^{1-\sigma}.
\]
The higher prime powers contribute $O_\sigma(1)$, as in the last
proof. This proves \eqref{eq:renewal-small-frequency}, and setting
$\tau=0$ proves \eqref{eq:renewal-zero-frequency}. Finally use
$\tau_R\to0$, continuity of $\psi$, and the asymptotic for $P_R$.
\end{proof}

Thus this signed operator supplies neither a universal contraction
in $\mathcal K_\sigma$ nor a uniform sign for its real Fourier
multiplier. Even a fixed nonvanishing smoothing factor near zero
does not remove its growing size. These are statements about a
generic operator estimate. They do not assert that the functions
testing its norm are Liouville sums, or that the particular $B$
belongs to $\mathcal K_\sigma$.
\footnote{The latter membership at every $\sigma>1/2$ is itself an
arithmetic assertion. It cannot be inserted as an unproved premise
in a proof of the critical energy bound.}
Moreover, the left side of \eqref{eq:prime-renewal-full} cannot be
replaced by $(\log X)B(X)$ without bounding their difference.
The renewal identity remains exact; the Liouville-specific upper
estimate remains missing.
